\documentclass[12pt, letterpaper]{article}

\usepackage[margin=0.4in]{geometry}
\usepackage{hyperref}
\usepackage{booktabs}
\usepackage{longtable}
\usepackage{parskip}
\usepackage{titlesec}
\usepackage{enumitem}

\titleformat{\section}{\large\bfseries}{}{0em}{}[\titlerule]
\titleformat{\subsection}{\normalsize\bfseries}{}{0em}{}

\hypersetup{
    colorlinks=true,
    linkcolor=black,
    urlcolor=blue,
    citecolor=black
}

\title{
    \textbf{Project 5} \\[0.4em]
    \large Architecture and Training Design Under a Hard Parameter Budget \\[0.4em]
    \normalsize DS 504: Big Data Analytics \textendash{} Spring 2026
}

\author{Shyam Patadia \\ \\ Worcester Polytechnic Institute}
\date{March 24, 2026}

\begin{document}

\maketitle
\thispagestyle{empty}

% -------------------------------------------------------
\section{Overview}
% -------------------------------------------------------

This project conducts a controlled ablation study of architectural
and training design choices for decoder-only language models
trained from scratch under a hard 16~MB artifact constraint,
as defined by the OpenAI Model Craft: Parameter Golf
Challenge~\cite{paramgolf2026}. The constraint requires that
model weights and training code combined not exceed 16,000,000
bytes. Model quality is evaluated by bits per byte (bpb) on
a held-out slice of the FineWeb~\cite{fineweb2024} validation
set, using the evaluation protocol from the official challenge
repository.

The experimental baseline is the NanoGPT-style decoder-only
transformer provided in the official challenge repository:
a 9-layer, 512-hidden-dimension model with a 1,024-token
SentencePiece vocabulary, tied input/output embeddings, and a
score of 1.2244~bpb~\cite{paramgolf2026}. All models are
initialized from scratch and trained on the FineWeb training
shard. No pretrained weights are used.

At FP32, the 16~MB budget accommodates approximately 4~million
parameters; at INT6, approximately 21~million. Quantization
precision therefore directly governs available model capacity,
rendering it a first-order design variable. The constraint
similarly makes positional encoding and optimizer state
footprint measurable in units of effective parameter count.
The existing literature on efficient language models typically
evaluates design choices in combination, confounding
attribution. This study varies each axis independently under
identical token budgets and per-configuration hyperparameter
tuning.

% -------------------------------------------------------
\section{Proposed Work}
% -------------------------------------------------------

\subsection{Model and Training Setup}

All configurations are decoder-only transformers trained from
scratch using \texttt{train\_gpt.py} from the official
Parameter Golf repository. The sp1024 SentencePiece vocabulary
is fixed across all experiments. Each configuration modifies
exactly one design axis relative to the baseline; parameters
freed by any modification are reallocated to model depth or
width such that total artifact size remains at or below 16~MB
across all conditions. All models are trained on WPI's
research cluster via SLURM.

\subsection{Research Objectives}

All primary configurations are evaluated on identical total
tokens from the same FineWeb training shard, trained to
convergence under early stopping on validation bpb. Learning
rate and warmup are tuned independently per configuration.
Results are reported as mean and standard deviation across
three independent random seeds. Exploratory conditions are
included where stable convergence is verified; results are
reported on a best-effort basis.

\begin{enumerate}[leftmargin=*, label=\textbf{O\arabic*.}]

    \item \textbf{Quantization Precision: BF16 vs.\ INT8
    vs.\ INT6.}
    Reducing precision from BF16 to INT8 doubles the effective
    parameter count within the artifact limit; INT6 triples it,
    making quantization the highest-leverage variable under
    this constraint. This objective examines whether the
    scaling laws for precision derived by Kumar et
    al.~\cite{scalingprecision} hold in the sub-21M parameter
    regime, and whether quantization noise from straight-through
    estimators during QAT erodes the capacity benefit. BF16
    and INT8 constitute the primary comparison; INT6 is
    included as an exploratory condition given its experimental
    support in \texttt{torchao}, with results reported where
    training converges stably. Bpb, training loss variance,
    and final artifact size are reported per precision level.

    \item \textbf{Positional Encoding: Learned vs.\ Full RoPE
    vs.\ None vs.\ Partial RoPE.}
    Learned absolute positional embeddings incur a parameter
    cost proportional to sequence length times hidden dimension.
    RoPE~\cite{rope} introduces no additional parameters,
    encoding relative position via
    $q_m^\top k_n = f(q,k,m-n)$. Freed parameters are
    reallocated to model capacity under a constant artifact
    budget. Learned embeddings, full RoPE, and no positional
    encoding constitute the primary conditions. Partial RoPE,
    which applies rotary embeddings to a subset of attention
    heads and requires direct attention kernel modification,
    is included as an exploratory condition with results
    reported where implementation is verified stable.

    \item \textbf{Optimizer: AdamW vs.\ Muon vs.\ Adafactor.}
    AdamW maintains per-parameter first and second moment
    estimates. Muon applies Nesterov momentum along the
    orthogonalized gradient. Adafactor factorizes the second
    moment, reducing optimizer-state memory from $O(n)$ to
    $O(\sqrt{n})$. Each optimizer is evaluated at its tuned
    learning rate optimum; reported metrics include
    convergence rate, final bpb, and total memory footprint
    including optimizer state.

\end{enumerate}

\subsection{Evaluation}

All configurations are evaluated on the held-out FineWeb
validation split from the official challenge repository,
ensuring reproducibility and comparability with existing
submissions. The 1.2244~bpb baseline constitutes the primary
reference. Results are reported as bpb delta over baseline
with standard deviation across seeds. All code, trained
checkpoints, and experiment logs will be released publicly
as a non-record submission to the official challenge
repository~\cite{paramgolf2026}.

% -------------------------------------------------------
\section{Technical Stack}
% -------------------------------------------------------

Python 3.11, PyTorch 2.x, HuggingFace \texttt{datasets} and
\texttt{tokenizers}, \texttt{torchao} for low-precision QAT,
Weights \& Biases for experiment tracking, SLURM for cluster
job scheduling.

% -------------------------------------------------------
\section{Schedule}
% -------------------------------------------------------

\begin{longtable}{@{}p{2.8cm} p{8cm} p{2.8cm}@{}}
\toprule
\textbf{Period} & \textbf{Tasks} & \textbf{Deliverable} \\
\midrule
\endfirsthead
Mar 24 -- Mar 30 &
Confirm baseline bpb; implement INT8
QAT via \texttt{torchao}; begin INT6
QAT integration; run LR sensitivity
sweep &
Baseline confirmed \\[4pt]
Mar 31 -- Apr 6 &
O1: Primary QAT runs (BF16, INT8)
to convergence; INT6 exploratory
runs where stable &
Quantization results \\[4pt]
Apr 7 -- Apr 13 &
O2: Primary positional encoding
runs (learned, full RoPE, none);
partial RoPE exploratory
implementation &
Positional encoding results \\[4pt]
Apr 14 -- Apr 20 &
O3: Optimizer comparison; progress
report draft &
Progress report draft \\[4pt]
\textbf{Apr 21} &
\textbf{Progress report due} &
\textbf{5-page report} \\[4pt]
Apr 22 -- Apr 28 &
Full analysis; convergence plots,
bpb tables, throughput results;
complete paper draft; presentation
slides; non-record submission &
Camera-ready report \\[4pt]
\textbf{May 5} &
\textbf{Final project due} &
\textbf{Report, code, slides} \\
\bottomrule
\end{longtable}

% -------------------------------------------------------
\section*{References}
\begingroup
\renewcommand{\section}[2]{}
\begin{thebibliography}{9}

\bibitem{paramgolf2026}
OpenAI.
\newblock {OpenAI Model Craft: Parameter Golf Challenge}, 2026.
\newblock \url{https://github.com/openai/parameter-golf}

\bibitem{fineweb2024}
Penedo, G. et al.
\newblock {The FineWeb Datasets: Decanting the Web for the Finest
Text Data at Scale}.
\newblock \textit{arXiv preprint arXiv:2406.17557}, 2024.

\bibitem{kaplan2020}
Kaplan, J. et al.
\newblock {Scaling Laws for Neural Language Models}.
\newblock \textit{arXiv preprint arXiv:2001.08361}, 2020.

\bibitem{hoffmann2022}
Hoffmann, J. et al.
\newblock {Training Compute-Optimal Large Language Models}.
\newblock \textit{arXiv preprint arXiv:2203.15556}, 2022.

\bibitem{scalingprecision}
Kumar, K. et al.
\newblock {Scaling Laws for Precision}.
\newblock \textit{ICLR}, 2025.
\newblock \textit{arXiv preprint arXiv:2403.08540}.

\bibitem{rope}
Su, J. et al.
\newblock {RoFormer: Enhanced Transformer with Rotary Position
Embedding}.
\newblock \textit{arXiv preprint arXiv:2104.09864}, 2021.

\end{thebibliography}
\endgroup

\end{document}